\section{The projector as a field}
\label{sec:fields}

Up to this point a projector has been a choice made once and for all. A choice of this kind becomes a field as soon as it is
allowed to vary from point to point, and the selection problem then changes character: instead of asking which projector is
right we ask which configurations an action prefers. Before any action can be written, three things must be known: the space in
which the projector moves, the connection it induces on its own image, and the sign of the natural metric on its variations. All
three can be computed exactly, and each of them carries a surprise.

\subsection{Configuration space}

Let $\mathcal P_r$ be the set of orthogonal projectors $P$ on $S_3=\R^{512}$ of rank $r$ that commute with the interior Lorentz
action $\Sigma$. Differentiating $P^2=P$ gives $dP\,P+P\,dP=dP$, hence
\[
P\,dP\,P=0,\qquad (\id-P)\,dP\,(\id-P)=0 :
\]
a variation of a projector is purely off-diagonal with respect to the splitting $\operatorname{im}P\oplus\ker P$. The tangent space
is therefore
\[
T_P\mathcal P_r=\bigl\{S:\ S^{\!\top}=S,\ PSP=0,\ (\id-P)S(\id-P)=0,\ [S,\Sigma]=0\bigr\},
\]
the tangent space of the real Grassmannian cut down by the Lorentz condition.

\begin{proposition}[Configuration space]
\label{prop:config}
The Lorentz commutant $\mathcal C\subset\End_\R(S_3)$ is spanned by the $32768$ blades that commute with the six interior
bivectors. Its centre is $\{\id,\Omega\}\cong\C$ with $\Omega=\gamma^0\gamma^1\gamma^2\gamma^3$, $\Omega^2=-\id$, and since a
semisimple real algebra of dimension $2\cdot128^2$ with centre $\C$ is $M_{128}(\C)$, we have $\mathcal C\cong M_{128}(\C)$.
Consequently
\[
\mathcal P_r\cong\mathrm{Gr}_\C\bigl(\tfrac r4,128\bigr),\qquad
\dim_\R\mathcal P_r=2\,\frac r4\Bigl(128-\frac r4\Bigr),
\]
which is $3584$ at $r=64$ and $6144$ at $r=128$, against $\dim\mathrm{Gr}_\R(64,512)=28672$ for the unconstrained real Grassmannian.
\end{proposition}

The complex structure that appears here is the interior volume $\Omega$, not the internal one $\omega_I$ of \cref{thm:i}. That is
not a contradiction. The commutant of the Lorentz action alone contains $e^{(1)}_2$, which commutes with the six interior bivectors
but anticommutes with the fourteen-index volume $\omega_I$, so $\omega_I$ is not central in $\mathcal C$; it becomes central only in
the smaller algebra that also commutes with the internal bivectors, which is the algebra of \cref{thm:i}. Each structure is central
in its own ambient algebra, which is the pattern described after \cref{prop:6plus8plus1}.

\begin{remark}[No blade is tangent]
\label{rem:noblade}
A count of tangent directions by blades fails, for a reason worth recording. Let $P=\prod_i\tfrac12(\id+X_i)$ be one of the
projectors of \cref{sec:projectors}, built from $k$ commuting involutions, and let $B$ be a blade anticommuting with at least one
$X_i$. Then $PBP=0$, but $(\id-P)B(\id-P)=(\id-P-P_\varepsilon)B$ with $P_\varepsilon=\prod_i\tfrac12(\id+\varepsilon_iX_i)$, which
vanishes only when $P+P_\varepsilon=\id$, that is only for $k=1$. There is also a rank argument: a blade is a signed permutation matrix,
invertible of rank $512$, while a tangent vector has the block form $\bigl(\begin{smallmatrix}0&A^{\!\top}\\A&0\end{smallmatrix}\bigr)$ and
rank at most $2r$. For $k\ge2$ no single blade is tangent, the tangent space is not spanned by blades, and a blade count
overestimates it, by a factor four at $r=64$ and two at $r=128$. The dimensions of \cref{prop:config} are obtained by imposing
both conditions sector by sector.
\end{remark}

\subsection{The sign of a variation}

The group that the geometry supplies is generated by the internal bivectors, and it moves a projector along an orbit inside
$\mathcal P_r$. Whether a bivector moves $P$ to another orthogonal projector, and with what sign the natural trace form measures
the motion, is decided by the signature rule of \cref{thm:sig}.

\begin{theorem}[The trace form splits by the Cartan involution]
\label{thm:ghost}
Let $P=P^{\!\top}=P^2$, let $B$ be a bivector and $T=[B,P]$. Then $T^{\!\top}=PB^{\!\top}-B^{\!\top}P$, and
\begin{enumerate}[label=(\roman*),leftmargin=2.2em]
\item if $B$ is compact, $B^{\!\top}=-B$, then $T$ is symmetric and $\Tr(T^2)=\Tr(T^{\!\top}T)>0$ unless $T=0$;
\item if $B$ is a boost, $B^{\!\top}=B$, then $T$ is antisymmetric and $\Tr(T^2)=-\Tr(T^{\!\top}T)<0$ unless $T=0$.
\end{enumerate}
Symmetric and antisymmetric matrices are trace-orthogonal, so on an orbit the trace form is an orthogonal sum of a positive
definite piece on the directions generated by compact bivectors and a negative definite piece on those generated by boosts. Only
the former are tangent to the manifold of orthogonal projectors.
\end{theorem}

\begin{proof}
$T^{\!\top}=(BP-PB)^{\!\top}=PB^{\!\top}-B^{\!\top}P$, which is $T$ when $B^{\!\top}=-B$ and $-T$ when $B^{\!\top}=B$. The trace of the
product of a symmetric and an antisymmetric matrix vanishes. A curve of orthogonal projectors has symmetric velocity, so the
antisymmetric directions leave that manifold.
\end{proof}

\begin{example}[The paired-mode projector at one constituent]
\label{ex:ghostPi}
Take the rank-four projector $\Pi$ of \cref{ex:pairedmode}. Its stabiliser in $\so{3,3}$ has dimension nine, so its orbit has
dimension six. On the bivectors
\[
B_I=\tfrac12\bigl(e_1e_2,\ e_1e_3,\ e_2e_3,\ e_1e_5,\ e_1e_6,\ e_2e_6\bigr),\qquad T_I=[B_I,\Pi],
\]
an exact $8\times8$ computation gives $\Tr(T_IT_J)=\operatorname{diag}(1,1,1,-1,-1,-1)$. The first three $B_I$ are time--time rotations,
the last three are boosts, exactly as \cref{thm:ghost} predicts. Only three of the fifteen basis bivectors commute with $\Pi$
individually, while the stabiliser has dimension nine, so six of its directions are combinations of basis bivectors and no single
one of them. Reading a stabiliser off basis elements overestimates the orbit.
\end{example}

Positive norm and tangency are nevertheless different conditions. For unit $T_{\rm sym}$ and $T_{\rm anti}$ of the two kinds,
$\Tr\bigl((T_{\rm sym}+\tfrac12T_{\rm anti})^2\bigr)=1-\tfrac14>0$, yet the direction is not symmetric and leaves the orthogonal
projectors. What is exact is the orthogonal splitting, not an equivalence between the sign of the norm and tangency. On the manifold
of orthogonal projectors itself every variation is symmetric and the trace form is positive definite, so the indefiniteness of the
example is an artefact of moving $\Pi$ along the full $\mathrm{Spin}(3,3)$ orbit. This is a statement about the target of a
projector field and not about a field theory over a base: in a base of split signature the kinetic contraction
$g^{\mu\nu}\Tr(D_\mu P\,D_\nu P)$ is indefinite in the spacetime index whatever the target metric is.

At three constituents the same reasoning counts the physical motions of $\Pfam$. Of the $91$ internal bivectors $43$ are compact,
forming the maximal compact $\so8\oplus\so6$, and the compact part of the stabiliser of $\Pfam$ has dimension nine
(\cref{thm:table}), so the compact orbit through $\Pfam$ has dimension
\[
43-9=34 ,
\]
which is also the rank of $B\mapsto[B,\Pfam]$ on the compact bivectors. Only seven of the $43$ compact basis bivectors commute with
$\Pfam$ individually, which repeats the warning of \cref{ex:ghostPi}. Of the $3584$ configuration directions at rank $64$,
thirty-four lie along the only group the geometry supplies. Whether these thirty-four are gauge redundancies or physical modes is a
question about an action, which the algebra does not supply.

\subsection{The induced connection over-generates}

A varying projector carries its own connection. If the gauge field is not postulated but induced, $\nabla=P\,d$ on the bundle of
images, then its curvature $F_{\mu\nu}=P[\partial_\mu P,\partial_\nu P]P$ is a composite, and the algebra in which it takes values is
computed rather than chosen. To leading order in the Ambrose--Singer theorem that algebra is
\[
\mathfrak h(P)=\operatorname{span}\bigl\{P[S,S']P:\ S,S'\in T_P\mathcal P_r\bigr\}.
\]

\begin{theorem}[Over-generation]
\label{thm:holonomy}
Let $P(x)$ range over all of $\mathcal P_{64}$. The curvature values $P[\partial_\mu P,\partial_\nu P]P$ are antisymmetric, hence
compact, and lie in the block, and their span on $\operatorname{im}\Pfam$ has dimension
\[
\dim\mathfrak h(\Pfam)=256=\dim\uu{16},
\]
the largest compact algebra the block admits at all. It exceeds the six dimensions of $\keff(\Pfam)=\uu1^3\oplus\su2$ by a factor
of about $43$.
\end{theorem}

The induced connection is not flat, every curvature value is compact, and every value lies in the block. The route fails by
over-generating and by nothing else. In particular chirality is not among the casualties.

\begin{proposition}[Chirality must be imposed, and costs nothing]
\label{prop:chiralpreserved}
Every projector of the family commutes with $\Gamma$. If a field $P(x)$ is required to satisfy $[P(x),\Gamma]=0$ everywhere, then
$[\partial_\mu P,\Gamma]=0$ and
\[
\bigl[P[\partial_\mu P,\partial_\nu P]P,\ \Gamma\bigr]=0 ,
\]
so the induced curvature preserves chirality and the holonomy lies in $\uu8\oplus\uu8$. That algebra is still far larger than
$\uu1^3\oplus\su2$.
\end{proposition}

The comparison of \cref{thm:holonomy} is between what the construction permits and what is wanted, not an impossibility theorem
about any particular field. A constant field has zero induced curvature, and a field compatible with $\Gamma$ has holonomy inside
$\uu8\oplus\uu8$. What the comparison does establish is the direction of the problem. A parent action must \emph{cut} the freedom of
the projector, by a constraint that shrinks $T_P\mathcal P$ before the curvature is formed or by a potential that freezes all but the
intended holonomy directions, rather than build a gauge field up from nothing. This reverses the usual motivation for composite
gauge fields. \Cref{fig:freedom} collects the numbers.

\begin{figure}[t]
\centering
\begin{tikzpicture}
\begin{axis}[xbar,xmode=log,log basis x=10,width=8.6cm,height=6.4cm,
  xmin=2,xmax=200000,
  ytick=data,
  symbolic y coords={keff,orbit,u8u8,u16,P64,GrR},
  yticklabels={{$\keff(\Pfam)=\uu1^3\oplus\su2$},{compact orbit of $\Pfam$},{curvature, chirality imposed},{curvature, unrestricted},{Lorentz-compatible projectors $\mathcal P_{64}$},{real Grassmannian $\mathrm{Gr}_\R(64,512)$}},
  yticklabel style={font=\small},
  xlabel={real dimension (logarithmic)},
  bar width=9pt,
  nodes near coords,
  nodes near coords style={font=\small,anchor=west},
  point meta=rawx,
  enlarge y limits=0.12,
  axis x line*=bottom, axis y line*=left]
\addplot[fill=tcol!45,draw=tcol] coordinates {(6,keff) (34,orbit) (128,u8u8) (256,u16) (3584,P64) (28672,GrR)};
\end{axis}
\end{tikzpicture}
\caption{How much freedom a projector field has at rank $64$, and how little of it is wanted. The configuration space of
Lorentz-compatible projectors (\cref{prop:config}) is far smaller than the real Grassmannian but still enormous; the group the
geometry supplies moves $\Pfam$ along only $34$ directions; and the curvature algebra the induced connection can reach
(\cref{thm:holonomy,prop:chiralpreserved}) exceeds the faithful gauge algebra of the block by more than an order of magnitude.}
\label{fig:freedom}
\end{figure}

\subsection{One explicit variational model}
\label{sec:trialaction}

A list of fields does not determine field equations. What can be done without a derived action is to adopt one explicitly, derive
its equation, and record what it can and cannot do. On a four-dimensional Lorentzian base, with $D_\mu P=\partial_\mu P+[A_\mu,P]$ and
with $f$ and $V$ as inputs, take
\begin{equation}
\label{eq:trialS}
S_P=\int\sqrt{-g}\,\Bigl(\tfrac{f^2}2\Tr\bigl(D_\mu P\,D^\mu P\bigr)-V(P)\Bigr).
\end{equation}
Varying along the Grassmannian, $\delta P=[\epsilon,P]$ with $\epsilon$ antisymmetric, gives
\begin{equation}
\label{eq:Pieom}
\bigl[P,\ f^2D_\mu D^\mu P+V'(P)\bigr]=0 ,
\end{equation}
and for $V=0$ this is the conservation law $D_\mu\mathcal C^\mu=0$ of the matrix current $\mathcal C_\mu=[P,D_\mu P]$. The variation
has to be projected onto the admissible tangent space whenever $P$ is restricted further, for instance to commute with $\Gamma$.

Three consequences are immediate. First, a potential that depends only on the spectrum of a projector of fixed rank is constant,
since $\Tr f(P)=rf(1)+(d-r)f(0)$; such a potential selects nothing, and for the pure kinetic term every constant admissible projector
is a critical point. Selection needs couplings that distinguish orbits, and there is more than one orbit to distinguish: at rank
$64$ the orbits through $\Pfam$ and $P_*$ have dimensions $34$ and $20$ (\cref{prop:orbits}). Second, whether varying $A_\mu$ gives a
Yang--Mills equation depends on how the action depends on $A$, not on $A$ appearing in a list of fields; a composite $A$ contributes
through the chain rule to the equations of its constituents. Third, the pointwise stabiliser condition $[A_\mu,P]=0$ is not the
condition $D_\mu P=0$ that a moving subbundle be preserved; the two agree only for constant $P$. Which connection acts on matter,
$P(d+A)$, an ambient connection preserving $P$, or an independent connection in a moving frame of the block, is a further choice,
and we make none of them here.

\subsection{What Yang--Mills theory keeps open}
\label{sec:ym}

Four conditions on a connection have been established as theorems: commutation with the Lorentz action
(\cref{thm:centraliser}); commutation with the projector, with a faithful action on its image (\cref{def:keff,thm:table});
antisymmetry, which for bivectors is \cref{thm:sig}; and preservation of the kinetic kernels $K^\mu$. They are homogeneous and
geometric, and together they determine the algebra in which $A_\mu$ takes its values, $\uu1^3\oplus\su2$ on the family block,
$\so6\oplus\su2^{\oplus3}$ at rank $256$ and the Pati--Salam algebra on $P_*$, over the single frame of \cref{prop:frame}. They do not
determine the field, since no parent action is derived; nor the couplings, since homogeneous conditions cancel every
normalisation; nor which abelian direction is electromagnetic, by \cref{thm:pairing}.

There is no representational obstruction behind this. By the theorem on universal connections \cite{NarasimhanRamanan}, every
connection on a bundle over a manifold of bounded dimension is the pullback of the canonical connection of a tautological bundle
over a Grassmannian of sufficient rank, so the projector representation of a gauge field is always available in principle. The
difficulty is dynamical, and power counting makes it concrete. A projector is dimensionless, so $\Tr(\partial P\,\partial P)$
carries a coupling of mass dimension $D-2$ and $\Tr F^2$ one of dimension $D-4$. In two dimensions the sigma model is
renormalisable and the composite gauge field acquires a kinetic term at leading order in $1/N$, with linear confinement of the
constituents \cite{DAddaDiVecchiaLuscher,Witten79}. In four dimensions the sigma model is not power-counting renormalisable and that
argument does not transfer.

Confinement is a separate question again. Invariant tensors permit meson-like and baryon-like operators, but they say nothing
about masses or about the energy of separated sources, and the diagnostics with dynamical matter are flux-tube and string-breaking
correlation matrices \cite{Bali05}. Moreover a block supplies a multiplicity index, not three constituents: three copies of a doublet
form a direct sum, whereas three constituents require a tensor product with antisymmetrisation, so even the operator statement needs
a Fock construction first.
